%============================================================
% CH2 · LESSON 5 — DFS AS A TOOLBOX
%============================================================
\subsection{DFS as a Toolbox}
\label{ssec:ch2-toolbox}

The four lessons so far looked at DFS from the inside. This lesson turns it around: here are the five things a contest solution most often asks DFS to do, each one a small change to the same skeleton. They are previews of the five patterns of Part~B.

\begin{diagrambox}
\begin{tikzpicture}[
    tool/.style={rectangle, rounded corners=3pt, draw=patternborder, fill=patternbg, line width=0.9pt,
                 font=\scriptsize, align=center, text width=2.45cm, minimum height=2.25cm},
    core/.style={rectangle, rounded corners=3pt, draw=deepblue, fill=theorybg, line width=1.1pt,
                 font=\scriptsize\ttfamily, align=left, text width=4.6cm}]
    \node[core] (c) at (0,0) {dfs(u):\\
\ \ \textcolor{successgreen}{\textbf{[enter u]}}\\
\ \ for v in adj[u]:\\
\ \ \ \ if v is white: dfs(v)\\
\ \ \ \ \textcolor{recursionred}{\textbf{[look at v's colour]}}\\
\ \ \textcolor{accentorange}{\textbf{[leave u]}}\\
\ \ \textcolor{softpurple}{\textbf{[un-mark u?]}}};
    \node[tool] (t1) at (-5.6,1.4) {\textbf{1. Cycles}\\(Pattern 1)\\[2pt]edge to a gray $v$ $\Rightarrow$ cycle; read it off the parents};
    \node[tool] (t2) at (-5.6,-1.4) {\textbf{2. Subtrees}\\(Pattern 2)\\[2pt]combine children at \emph{leave}; pass values down at \emph{enter}};
    \node[tool] (t3) at (5.6,1.4) {\textbf{3. Timestamps}\\(Pattern 3)\\[2pt]$\tin{}$ at enter, $\tout{}$ at leave; subtrees are intervals};
    \node[tool] (t4) at (5.6,-1.4) {\textbf{4. Backtracking}\\(Pattern 4)\\[2pt]un-mark at leave $\Rightarrow$ every simple path};
    \node[tool] (t5) at (0,-2.9) {\textbf{5. Reverse search}\\(Pattern 5)\\[2pt]run it on reversed edges, from the target};
    \draw[flow] (t1.east) -- (c.west |- t1.east);
    \draw[flow] (t2.east) -- (c.west |- t2.east);
    \draw[flow] (t3.west) -- (c.east |- t3.west);
    \draw[flow] (t4.west) -- (c.east |- t4.west);
    \draw[flow] (t5.north) -- (c.south);
\end{tikzpicture}
\figcap{One skeleton, five tools. Each tool fills in one of the coloured slots.}
\end{diagrambox}

%------------------------------------------------------------
\subsubsection*{Tool 1 and 2: information up and down a rooted tree}

Most ``subtree'' problems give a \textbf{tree} as $n-1$ undirected edges and fix a root. DFS from the root \emph{is} the rooted tree: every edge is a tree edge (there are no back edges in a tree), and the only neighbour that is not a child is the parent. So instead of a \texttt{visited} array one usually passes the parent:

\begin{lstlisting}[style=cpp]
#include <bits/stdc++.h>
using namespace std;

// A rooted tree from undirected edges: depth goes DOWN (enter), size goes UP (leave).
vector<vector<int>> adj;
vector<int> depth_, sz;

void dfs(int u, int p) {                 // p = parent of u (0 for the root)
    depth_[u] = (p == 0) ? 0 : depth_[p] + 1;   // pre-order: the parent is already done
    sz[u] = 1;
    for (int v : adj[u])
        if (v != p) {                    // in a tree, "not the parent" == "not visited"
            dfs(v, u);
            sz[u] += sz[v];              // post-order: the child is finished
        }
}

int main() {
    int n;
    cin >> n;
    adj.assign(n + 1, vector<int>());
    for (int i = 0; i < n - 1; i++) {
        int a, b;
        cin >> a >> b;
        adj[a].push_back(b);
        adj[b].push_back(a);
    }
    depth_.assign(n + 1, 0); sz.assign(n + 1, 0);
    dfs(1, 0);
    for (int v = 1; v <= n; v++) cout << v << ": depth " << depth_[v] << ", size " << sz[v] << "\n";
}
\end{lstlisting}

\begin{diagrambox}
\begin{tikzpicture}[x=1.1cm, y=1cm]
    \node[gvsrc] (r1) at (0,0) {1};
    \node[gv] (r2) at (-1.8,-1) {2}; \node[gv] (r3) at (0,-1) {3}; \node[gv] (r4) at (1.8,-1) {4};
    \node[gv] (r5) at (-2.5,-2) {5}; \node[gv] (r6) at (-1.1,-2) {6}; \node[gv] (r7) at (1.8,-2) {7};
    \node[gv] (r8) at (1.8,-3) {8};
    \foreach \a/\b in {1/2, 1/3, 1/4, 2/5, 2/6, 4/7, 7/8} \draw[gtree] (r\a) -- (r\b);
    \foreach \v/\d/\s in {1/0/8, 2/1/3, 3/1/1, 4/1/3, 5/2/1, 6/2/1, 7/2/2, 8/3/1}{
        \node[gdist, text=accentorange, above left=-1pt and -1pt of r\v] {\tiny d\d};
        \node[gdist, text=primaryblue, above right=-1pt and -1pt of r\v] {\tiny s\s};
    }
    \draw[->, thick, accentorange] (3.4,0) -- (3.4,-3) node[midway, right, font=\scriptsize, align=left] {depth flows\\down (enter)};
    \draw[->, thick, primaryblue] (5.3,-3) -- (5.3,0) node[midway, right, font=\scriptsize, align=left] {size flows\\up (leave)};
\end{tikzpicture}
\figcap{The tree $T$ with edges $1{-}2, 1{-}3, 1{-}4, 2{-}5, 2{-}6, 4{-}7, 7{-}8$, rooted at 1: orange $d$ = depth, blue $s$ = subtree size, both printed by the program.}
\end{diagrambox}

\begin{lstlisting}[style=trace]
input : 8 / 1 2 / 1 3 / 1 4 / 2 5 / 2 6 / 4 7 / 7 8
output: 1: depth 0, size 8     2: depth 1, size 3     3: depth 1, size 1     4: depth 1, size 3
        5: depth 2, size 1     6: depth 2, size 1     7: depth 2, size 2     8: depth 3, size 1
\end{lstlisting}

Every Pattern~2 problem is this program with a different quantity in place of the size (a count of colours, a sum, a product of sizes), or in place of the depth (a sum of values along the root path).

%------------------------------------------------------------
\subsubsection*{Tool 3: path reconstruction}

DFS leaves behind \texttt{parent[]}. Walking it from $t$ back to $s$ gives \emph{a} path from $s$ to $t$:
\begin{lstlisting}[style=snippet]
vector<int> path;
for (int x = t; x != s; x = parent[x]) path.push_back(x);
path.push_back(s);
reverse(path.begin(), path.end());     // s ... t
\end{lstlisting}
On $G_3$ from $s = 1$ the parents are $\texttt{parent}[2] = 1$, $\texttt{parent}[3] = 2$, $\texttt{parent}[5] = 2$, $\texttt{parent}[4] = 5$, $\texttt{parent}[6] = 5$, and the path to $4$ is $1, 2, 5, 4$ --- three edges, although the edge $1$--$4$ exists. DFS paths are paths, not shortest paths; for shortest paths use BFS (Chapter~3).

%------------------------------------------------------------
\subsubsection*{Tool 4: backtracking --- un-mark on the way out}

DFS marks a vertex once and never visits it again; that is what makes it linear. Remove the mark when the vertex is \emph{left} and the search changes completely: it now explores every \textbf{simple path} from the start --- every way of walking without repeating a vertex. The marks then mean ``on the current path'', exactly the gray vertices.

\begin{lstlisting}[style=cpp]
#include <bits/stdc++.h>
using namespace std;

// Backtracking: print EVERY simple path from s to t. onPath[v] = v is on the current path.
vector<vector<int>> adj;
vector<bool> onPath;
vector<int> path;
int t, found = 0;

void go(int u) {
    onPath[u] = true; path.push_back(u);
    if (u == t) {
        found++;
        for (int x : path) cout << x << " ";
        cout << "\n";
    } else {
        for (int v : adj[u])
            if (!onPath[v]) go(v);
    }
    onPath[u] = false; path.pop_back();      // un-mark: other paths may use u again
}

int main() {
    int n, m, s;
    cin >> n >> m;
    adj.assign(n + 1, vector<int>());
    for (int i = 0; i < m; i++) {
        int a, b;
        cin >> a >> b;
        adj[a].push_back(b);
        adj[b].push_back(a);
    }
    cin >> s >> t;
    onPath.assign(n + 1, false);
    go(s);
    cout << found << " simple paths\n";
}
\end{lstlisting}

\begin{lstlisting}[style=trace]
input : G3, then s = 1, t = 6
output: 1 2 5 6
        1 4 5 6
        2 simple paths
\end{lstlisting}

The price: the number of simple paths can be exponential (a $6 \times 6$ grid graph already has $1\,262\,816$ simple paths between opposite corners --- the backtracking program counted them), so backtracking is only for tiny graphs or heavily pruned searches. Pattern~4 is about recognising when it is affordable.

%------------------------------------------------------------
\subsubsection*{Tool 5: reverse search}

``From which vertices can $t$ be reached?'' would need one search per start vertex. Reverse every edge instead: $u$ can reach $t$ in $G$ iff $t$ can reach $u$ in the reversed graph $G^R$. One search from $t$ in $G^R$ answers the question for all $u$ at once.
\begin{lstlisting}[style=snippet]
for (auto& e : edges) radj[e.second].push_back(e.first);   // edge a -> b becomes b -> a
dfs_on(radj, t);                                            // marked = everything that can reach t
\end{lstlisting}
On $D$ the vertices that can reach 4 are $\{1, 2, 3, 4, 5\}$: 3 gets there via $3 \to 1 \to 2 \to 4$, while 6, 7 and 8 cannot (the only edges out of $\{6, 7\}$ stay inside it). Pattern~5 also reverses \emph{roles}: start from the exit of a maze, from the border of a grid, from the sinks of a DAG.

%------------------------------------------------------------
\subsubsection*{Which tool? A first decision table}

\begin{center}
\small
\renewcommand{\arraystretch}{1.3}
\begin{tabularx}{\linewidth}{@{}>{\raggedright\arraybackslash}X l l@{}}
\toprule
\textbf{The question is about \dots} & \textbf{tool} & \textbf{cost} \\
\midrule
``is there a cycle / a deadlock / a contradiction in the order?'' & 1: colours, back edges & \VE \\
``for every vertex, something about the vertices below it'' & 2: post-order aggregation & \VE \\
``for every vertex, something about the path from the root'' & 2: pre-order push-down & \VE \\
``is $u$ an ancestor of $v$?'', ``the $k$-th vertex of $u$'s subtree in DFS order'' & 3: timestamps & \VE{} $+\ \bigO{1}$ per query \\
``all paths'', ``the best path'' on $\le 15$--$25$ vertices or cells & 4: backtracking & exponential \\
``which vertices can reach the target / the border / the exit?'' & 5: reverse search & \VE \\
\bottomrule
\end{tabularx}
\end{center}

\subsubsection*{Check your understanding}

\begin{enumerate}[leftmargin=*]
    \item In the tree $T$, re-root at 4. What are the new depth and subtree size of vertex 1?
    \item Why does the tree DFS need no \texttt{visited} array, while a DFS on $G_3$ does?
    \item The backtracking program is run on $G_3$ with $s = 3$, $t = 6$. How many simple paths does it print?
    \item Which vertices of $D$ can reach vertex 7?
    \item Without backtracking, can the plain DFS (marks never removed) print \emph{all} simple paths from $s$ to $t$? Give an example on $G_3$.
\end{enumerate}

\begin{solutionbox}[Answers]
\begin{enumerate}[leftmargin=*]
    \item Rooted at 4, vertex 1 is a child of 4: depth 1. Its subtree is everything except the branch $\{4, 7, 8\}$: size $8 - 3 = 5$. (The ``$n - \mathrm{sz}$'' trick returns in Problem~\ref{ssec:ch2-highest-score}.)
    \item In a tree the only already-visited neighbour of $u$ is its parent; skipping the parent is enough. $G_3$ has a cycle $1, 2, 5, 4$, so a vertex can be reached a second time along a non-tree edge.
    \item Every path from 3 must start $3, 2$; from 2 it continues $2, 5, 6$ or $2, 1, 4, 5, 6$: \textbf{2} paths (the program prints \texttt{3 2 1 4 5 6}, then \texttt{3 2 5 6}, because \texttt{adj[2]} lists 1 before 5).
    \item Search from 7 in the reversed graph: $7 \leftarrow 6$, $7 \leftarrow 8$; $6 \leftarrow 5$, $6 \leftarrow 1$; $1 \leftarrow 3$; $3 \leftarrow 2$, $3 \leftarrow 4$. So \textbf{all 8 vertices} can reach 7 --- even 4, along $4 \to 3 \to 1 \to 6 \to 7$.
    \item No: once 5 is marked by the path $1, 2, 5, 6$, the path $1, 4, 5, 6$ can never use 5 again. Plain DFS finds one path per reachable vertex.
\end{enumerate}
\end{solutionbox}
